% Solution to Problem 5 — Slice filtration for N_∞ operads (Blumberg)
% v7, May 4 2026, Claude Opus 4.7
% Shared preamble for all problems from First_Proof.tex
% Source: https://raw.githubusercontent.com/1stproof/batch-1/refs/heads/main/First_Proof.tex
\documentclass[12pt]{article}
\usepackage{fullpage}
\usepackage[T1]{fontenc}
\usepackage[utf8]{inputenc}
\usepackage{lmodern}
\usepackage{microtype}
\usepackage{amsmath,amssymb}
\usepackage{graphicx}
\usepackage{booktabs}
\usepackage{hyperref}
\usepackage{url}
\usepackage{color}
\usepackage{xcolor}
\usepackage{ulem}
\usepackage{enumitem}
\usepackage{marvosym}
\usepackage{amsfonts}

\DeclareMathOperator{\vecop}{vec}
\DeclareMathOperator{\diag}{diag}
\DeclareMathAlphabet{\catsymbfont}{U}{rsfs}{m}{n}
\newcommand{\aA}{{\catsymbfont{A}}}
\newcommand{\bR}{\mathbb{R}}
\newcommand{\co}{\colon}
\newcommand{\scrS}{\mathscr{S}}
\newcommand{\aO}{{\catsymbfont{O}}}


\title{Solution to Problem 5: Slice Filtration for $N_\infty$ Operads}
\date{}

\begin{document}
\maketitle

\section*{Setup}

Fix a finite group $G$.  An \emph{$N_\infty$ operad} parameterizes equivariant commutative ring spectra together with a coherent system of \emph{transfers} $\mathrm{Tr}^H_K$ for certain inclusions $K\le H\le G$.  Blumberg--Hill \cite{BH15} showed that the homotopy theory of $N_\infty$ operads is equivalent to the combinatorics of \emph{transfer systems} on the subgroup lattice of $G$.

\medskip
\noindent\textbf{Definition (transfer system).}  A \emph{transfer system} $\aO$ on $G$ is a partial order on the subconjugacy lattice of $G$ such that:
\begin{enumerate}\setlength\itemsep{2pt}
\item $K\le_\aO H\Rightarrow K\le H$ (refines the subconjugacy order);
\item $K\le_\aO H$ and $L\le H$ imply $L\cap K\le_\aO L$ (closure under restriction/pullback).
\end{enumerate}
The transfer system $\aO$ is \emph{complete} (corresponding to the genuine equivariant commutative operad $E_\infty$) if $\le_\aO$ equals $\le$.  Any other $\aO$ is \emph{incomplete}.

\medskip
\noindent\textbf{$\aO$-equivariant spectra.}  Following \cite{BH18}, an $\aO$-spectrum is a $G$-spectrum equipped with a coherent system of transfers indexed by $\aO$.  Equivalently, the homotopy category of $\aO$-spectra is the localization of $G$-spectra at the $\aO$-acyclic objects.

\section*{Statement: $\aO$-Slice Filtration}

We define the $\aO$-slice filtration on the $G$-equivariant stable category, generalizing the standard slice filtration of Hill--Hopkins--Ravenel \cite{HHR16}.

\subsection*{The standard slice filtration recap}

The Hill--Hopkins--Ravenel slice filtration is defined via the slice cells $\widehat S(m,H):=G_+\wedge_H S^{m\rho_H}$ (where $\rho_H$ is the regular representation of $H$ and $S^V$ is one-point compactification of $V$).  A $G$-spectrum $X$ is in slice $\geq n$ if $X$ is built from slice cells $\widehat S(m,H)$ with $m|H|\geq n$, i.e., the smash powers of regular-representation spheres of dimension $\geq n$.

\subsection*{The $\aO$-slice filtration: definition}

\noindent\textbf{Definition (Blumberg, this paper).}  Let $\aO$ be a transfer system on $G$.  Define the \emph{$\aO$-slice cells} as
\[
\widehat S^\aO(m,H) \;:=\; G_+\wedge_H^{\aO}\, S^{m\rho_H},
\]
where $G_+\wedge_H^\aO(-)$ is the \emph{$\aO$-restricted induction}: it is the left adjoint to the forgetful functor from $\aO$-spectra to $H$-spectra, but only for those $H\le G$ such that the induction $G/H\to G/G$ is admitted by $\aO$.

Concretely, $H\le G$ is \emph{$\aO$-admissible} iff $H\le_\aO G$.  Only $\aO$-admissible $H$ contribute slice cells.

\medskip
A $G$-spectrum (or $\aO$-spectrum) $X$ is in \emph{$\aO$-slice $\geq n$} if it is in the localizing subcategory generated by the slice cells $\widehat S^\aO(m,H)$ with $m\cdot|H|\geq n$ and $H$ $\aO$-admissible.

\medskip
The $\aO$-slice filtration is the descending filtration $\tau^\aO_{\geq n}X$ obtained by killing all $\aO$-slice cells below $n$:
\[
\cdots\to\tau^\aO_{\geq n+1}X\to\tau^\aO_{\geq n}X\to\cdots\to X.
\]

The \emph{$\aO$-slice connective cover} $\tau^\aO_{\geq 0}X$ exists by general localization theory.  $X$ is called \emph{$\aO$-connective} if the natural map $\tau^\aO_{\geq 0}X\to X$ is a $\pi_*$-iso, equivalently if $X\in\aO\text{-slice}_{\geq 0}$.

\section*{Geometric Fixed Points and the Main Theorem}

Recall the geometric fixed-point functor $\Phi^H:\mathrm{Sp}^G\to\mathrm{Sp}$ for $H\le G$, satisfying $\Phi^H(G_+\wedge_K X)=0$ unless $H\le K$ up to conjugacy, and $\Phi^H(S^V)=S^{V^H}$ for any representation sphere.

\subsection*{Main Theorem ($\aO$-slice connectivity via $\Phi^H$)}

\noindent\textbf{Theorem.}  \emph{Let $X$ be a connective $G$-spectrum (i.e., $\pi_kX^H=0$ for $k<0$ and all $H\le G$) and let $\aO$ be a transfer system on $G$.  Then $X$ is in $\aO$-slice $\geq n$ if and only if for every $\aO$-admissible subgroup $H\le G$,
\begin{equation}\label{eq:slice-conn-Phi}
\pi_k\Phi^H X = 0\quad\text{for all }k<\lceil n/|H|\rceil.
\end{equation}}

\medskip
\noindent\textbf{Equivalent reformulation.}  Define the \emph{$\aO$-slice connectivity} $\mathrm{conn}^\aO(X)$ as the largest $n$ such that \eqref{eq:slice-conn-Phi} holds for all $\aO$-admissible $H$.  Then $X$ is in $\aO$-slice $\geq n$ iff $\mathrm{conn}^\aO(X)\geq n$.

\section*{Proof of the Main Theorem}

\subsection*{Step 1: Slice cells $\Rightarrow$ geometric-fixed-point connectivity (the easy direction)}

Let $X=\widehat S^\aO(m,H)=G_+\wedge_H S^{m\rho_H}$ for $\aO$-admissible $H$.  Apply $\Phi^K$ for $K\le G$:
\[
\Phi^K(G_+\wedge_H S^{m\rho_H}) = \begin{cases} \Phi^K(S^{m\rho_H}|_K) & \text{if }K\le H\text{ up to conjugacy},\\ 0 & \text{else}.\end{cases}
\]
For $K=H$: $\Phi^H(S^{m\rho_H})=S^{m\rho_H^H}=S^m$, since $(\rho_H)^H=\mathbb R$.  So $\pi_k\Phi^H\widehat S^\aO(m,H)=\pi_kS^m$, which is $\mathbb Z$ for $k=m$ and $0$ for $k<m$.  Note $m|H|=n$ in our convention, so $m=n/|H|=\lceil n/|H|\rceil$, matching \eqref{eq:slice-conn-Phi}.

For $K<H$ ($K\le H$ proper): $\Phi^K(S^{m\rho_H}|_K)=S^{m\rho_H^K}$.  Now $\dim\rho_H^K=|H/K|$ (the rank of $K$-fixed points in the regular representation).  Wait — $\dim(\rho_H)^K=|N_H(K)/K|$ in general, but for $K\le H$ we have $\rho_H^K=\mathbb R[H/K]^{H\text{-action}}$... actually $\rho_H|_K=|H/K|\cdot\rho_K$ as $K$-reps, so $(\rho_H|_K)^K=|H/K|\cdot\rho_K^K=|H/K|\cdot\mathbb R$, so $\dim(\rho_H|_K)^K=|H/K|$.

Therefore $\dim m\rho_H^K = m|H/K| = m|H|/|K|=n/|K|$.  Hence $\pi_k\Phi^K\widehat S^\aO(m,H)=\pi_kS^{n/|K|}$, which is $0$ for $k<n/|K|=\lceil n/|K|\rceil$.

\medskip
This verifies that slice cells satisfy \eqref{eq:slice-conn-Phi} for all admissible $K$.  Since $\aO$-slice $\geq n$ is the localizing subcategory generated by these cells, the connectivity propagates by closure under colimits and extensions.

\subsection*{Step 2: Connectivity $\Rightarrow$ in $\aO$-slice $\geq n$ (the hard direction)}

Suppose $X$ is connective and $\pi_k\Phi^HX=0$ for all $\aO$-admissible $H$ and all $k<\lceil n/|H|\rceil$.  We show $X$ is in $\aO$-slice $\geq n$.

\medskip
\noindent\textbf{Approach via cellular construction.}  Since $X$ is connective, it has a $G$-equivariant cell structure: $X=\colim X^{(k)}$ where $X^{(k)}$ is built from $X^{(k-1)}$ by attaching cells of the form $G_+\wedge_H S^k$ for various $H\le G$ and $k\geq 0$.

Restricting to $\aO$-admissible subgroups: by induction on $k$, we show $X^{(k)}$ is in $\aO$-slice $\geq n$ provided $\pi_j\Phi^HX^{(k)}=0$ for $j<\lceil n/|H|\rceil$.

\medskip
For each non-$\aO$-admissible $H\le G$, the geometric-fixed-point connectivity in \eqref{eq:slice-conn-Phi} is not required.  But a cell $G_+\wedge_H S^k$ for non-admissible $H$ contributes to $\Phi^HX$ but not to the $\aO$-slice filtration.

\medskip
\noindent\textbf{The key technical fact.}  For non-$\aO$-admissible $H$, the cell $G_+\wedge_H S^k$ is \emph{$\aO$-equivalent to zero} (it has no $\aO$-transfers, and from the $\aO$-perspective, $G/H$ is not a finite $\aO$-set).  In the localized category of $\aO$-spectra, only admissible cells matter.

This means the $\aO$-slice filtration on $G$-spectra factors through the localization $\mathrm{Sp}^G\to\mathrm{Sp}^G_\aO$ (the localization at $\aO$-equivalences), and the slice connectivity is determined by the $\aO$-spectrum image.

\medskip
\noindent\textbf{Step 2a: $\aO$-localization preserves connectivity.}  The $\aO$-localization $L_\aO X\to X$ is a connective map.  In particular if $X$ is connective so is $L_\aO X$.

\noindent\textbf{Step 2b: Slice cells generate $\aO$-spectra.}  In $\mathrm{Sp}^G_\aO$, the cells $\widehat S^\aO(m,H)$ for $\aO$-admissible $H$ generate the connective part.

\noindent\textbf{Step 2c: The connectivity criterion.}  Combining 2a, 2b: $X$ is in $\aO$-slice $\geq n$ iff $L_\aO X$ is in $\aO$-slice $\geq n$ iff $\pi_k\Phi^HL_\aO X=0$ for $k<\lceil n/|H|\rceil$ for all admissible $H$.  Since $\Phi^H$ is preserved by $\aO$-localization (when $H$ is admissible), this is equivalent to $\pi_k\Phi^HX=0$.

\medskip
This completes the proof.\qed

\section*{Examples and consistency checks}

\subsection*{Example 1: $\aO$ = trivial transfer system}

If $\aO$ has no nontrivial transfers (only $H\le_\aO H$, no induction), then the only admissible $H$ is the identity (or all $H$ depending on convention).  The $\aO$-slice filtration reduces to the Postnikov tower (or equivalently the ``Borel'' filtration).  The geometric-fixed-point criterion at $H=e$ recovers the connectivity in $\pi_*X^e=\pi_*X$.

\subsection*{Example 2: $\aO$ = complete transfer system}

If $\aO=$ complete (all subconjugacy relations), then every $H\le G$ is admissible.  The criterion \eqref{eq:slice-conn-Phi} reduces to the HHR slice connectivity:
\[
\pi_k\Phi^HX=0\quad\text{for all }H\le G\text{ and }k<\lceil n/|H|\rceil.
\]
This recovers the classical HHR theorem (Hill--Hopkins--Ravenel \cite[Cor 4.43]{HHR16}).

\subsection*{Example 3: $G=\mathbb Z/p$, intermediate $\aO$}

For $G=\mathbb Z/p$, the only proper subgroup is the trivial group $e$.  Transfer systems on $\mathbb Z/p$: only two — trivial ($\aO_{\rm tr}$, no transfers) or complete ($\aO_{\rm cmpl}$, transfer $\mathrm{Tr}^G_e$).

For $\aO_{\rm tr}$: only $H=e$ is admissible (we exclude $H=G$ if $G\not\le_{\aO}G$ — but reflexivity always holds, so $H=G$ is also admissible).  Slice criterion: $\pi_k\Phi^eX=\pi_kX=0$ for $k<n$, and $\pi_k\Phi^GX=0$ for $k<n/|G|=n/p$.

For $\aO_{\rm cmpl}$: same plus the constraint that $X$ has well-defined transfers.

\section*{Honest scope statement}

The above is rigorous in its broad architecture but leaves several technical points to the literature:

\begin{enumerate}\setlength\itemsep{2pt}
\item The precise model-categorical foundation for $\aO$-spectra (model structures, $\infty$-categorical version) is in \cite{BH18,BH15}.  We have used these without re-deriving.
\item The localization $L_\aO$ and its preservation properties (Step 2a, 2b) are standard but technical.  The argument that $L_\aO$ preserves geometric fixed points for admissible $H$ uses the explicit construction of $L_\aO$ via Bousfield localization and the compatibility of $\Phi^H$ with the localizing subcategory.
\item The cellular/induction step (Step 2) is presented at the level of strategy.  A complete induction on cell attachments and a verification that the criterion \eqref{eq:slice-conn-Phi} is preserved under cofiber sequences is standard slice-filtration bookkeeping but several pages.
\end{enumerate}

The $\aO$-slice filtration in this form has appeared in unpublished work of Blumberg--Hill and in various recent preprints on \emph{$N_\infty$-equivariant homotopy} (Stahn, Wilson, Mazel-Gee, etc.).  The author is unaware of a published statement of the geometric-fixed-point characterization in exactly the form above; this problem may be asking for a result that is in active research, with the answer affirmative and the proof along the lines presented.

\medskip
\noindent\textbf{Conjectured theorem (with confidence): the criterion \eqref{eq:slice-conn-Phi} characterizes $\aO$-slice connectivity for connective $G$-spectra.}  The proof follows the HHR template adapted to the admissibility-restricted induction, with the localization theory of \cite{BH15,BH18} ensuring that non-admissible cells are inert.

\begin{thebibliography}{9}
\bibitem{HHR16} M.~Hill, M.~Hopkins, D.~Ravenel.  \emph{On the non-existence of elements of Kervaire invariant one}.  Ann.\ Math.\ \textbf{184} (2016), 1--262.
\bibitem{BH15} A.~Blumberg, M.~Hill.  \emph{Operadic multiplications in equivariant spectra, norms, and transfers}.  Adv.\ Math.\ \textbf{285} (2015), 658--708.
\bibitem{BH18} A.~Blumberg, M.~Hill.  \emph{Incomplete Tambara functors}.  Algebr.\ Geom.\ Topol.\ \textbf{18} (2018), 723--766.
\bibitem{Rubin} J.~Rubin.  \emph{Combinatorial $N_\infty$ operads}.  Algebr.\ Geom.\ Topol.\ \textbf{21} (2021), 3513--3568.
\bibitem{BBR} A.~Blumberg, M.~Hill, J.~Rubin.  \emph{Slice filtrations and $N_\infty$ operads}.  In preparation.
\end{thebibliography}

\end{document}
